\chapter{Choosing a segment, reading a posting}
\label{ch:segment}

\section*{What this chapter is for}

Before you write a line of an application, two decisions narrow everything that
follows: which kind of company you are aiming at, and what the specific posting
in front of you is actually going to score. Both are cheap. Neither is usually
made deliberately.

\section{Three segments, wanting three different people}

The most useful piece of targeting advice found while assembling this book came
from an engineer who interviewed at a few dozen companies in one stretch and
concluded the market divides into three: \reported[single]{companies building
products on top of models, companies building the infrastructure that serves
them, and companies building the models themselves \dash{} and that you should
pick one.}{kalra-transition}

They want genuinely different things.

\begin{kittable}{@{}lXX@{}}
\textbf{Segment} & \textbf{What they are buying} & \textbf{What they will test} \\
\midrule
Product & judgement about users, speed to a working thing, taste about when the model is the wrong tool & discovery, a build round, product sense \\
Infrastructure & distributed systems, reliability, cost control, APIs consumed by people you cannot talk to & system design, failure modes, throughput and tail latency \\
Model & research proximity, evaluation depth, comfort with the parts of the stack this book calls noise & the parts of the stack this book calls noise \\
\end{kittable}

\judgement{For an engineer arriving from backend or platform work, the
infrastructure segment is usually the strongest fit and the least applied to.
Its postings often use product vocabulary, and its interviews turn out to be
the distributed-systems interviews you have been passing for years with one
unfamiliar dependency in the middle.}

Picking a segment is not a life decision. It decides which artefacts you build
in Part~\ref{part:evidence}, which of the bank's chapters you study first, and
which half of your existing experience you lead with.

\section{Reading a posting as a rubric}

A posting is a marketing document with a rubric inside it. The rubric is
recoverable, and recovering it takes about twenty minutes.

\begin{kitmethod}
\textbf{1. Separate the three sections.} Almost every posting has responsibilities
(what you will do), requirements (what they will check), and culture (what they
will check without telling you). They are scored by different people.

\textbf{2. Extract the verbs from responsibilities.} These map to interview
rounds more reliably than anything else on the page. \emph{Design}, \emph{own},
\emph{define} imply a design round and a scope conversation. \emph{Integrate},
\emph{ship} imply a build round. \emph{Mentor}, \emph{raise the bar},
\emph{influence} imply a behavioural round with the Staff questions in it.

\textbf{3. Count the repeats.} A word that appears three times in one posting is
a word somebody argued about. It will appear again in the interview, and using
it back is not flattery, it is speaking the language the rubric was written in.

\textbf{4. Find the level definition.} If the company publishes a ladder, the
posting's band maps to a row in it. That row is the rubric, stated by them,
which is better than any inference you can make from the advert.

\textbf{5. Write the rubric down as rows.} Each responsibility becomes a row:
what they said, which round will test it, what evidence you have, and how
strong that evidence is. \code{templates/role-rubric.md} is that table.
\end{kitmethod}

The output is a page that tells you where you are weak \emph{before} the first
call, which is the only point at which the information is still actionable.

\section{The things a posting tells you that are not about the work}

\textbf{The band, where it is published.} Some jurisdictions require a range.
Where you have one, it is a fact about the level and not an invitation, and
Chapter~\ref{ch:company} covers what to do with it.

\textbf{The location and work model, stated in the posting's own metadata
rather than its prose.} Prose says \emph{flexible}. The metadata field says
\emph{on-site}. Chapter~\ref{ch:fit} is about taking the metadata seriously
before spending a week.

\textbf{Whether the role is advertised in several cities at once.} That usually
means a team assignment rather than a remote role, whatever the body text says.

\textbf{Whether evaluation, observability or cost appear anywhere.} Their
presence means somebody has operated one of these systems and been hurt.
\judgement{It is the strongest single signal in a posting that the team is past
the demo stage, and it is worth more than the stack list.}

\begin{drill}
Take the posting you are most serious about. Fill in
\code{templates/role-rubric.md} completely, including the evidence-strength
column, and do not skip a row because you know the answer. Then count the rows
where your evidence is weak. If it is most of them, the useful conclusion is
not that you should prepare harder; it is that you should read
Chapter~\ref{ch:fit} before you apply.
\end{drill}

\section*{What this chapter does not claim}

The three-segment split rests on one engineer's account of one stretch of
interviewing, which is why the sentence that introduces it says so. It is used
here because it is a useful way to organise a decision you have to make anyway,
not because it has been measured.

This chapter also does not claim that postings are written carefully. Many are
assembled from a template and a previous role, and a requirement can be there
because nobody deleted it. That is a reason to weight repeats and verbs over
individual bullet points \dash{} not a reason to ignore the document.
